\documentclass[UTF8,11pt,a4paper,fontset=windows]{ctexart}

% ============================================================
%  Packages
% ============================================================
\usepackage{geometry}
\geometry{left=2.3cm,right=2.3cm,top=2.7cm,bottom=2.7cm,headheight=22.1pt,headsep=0.55cm,footskip=0.9cm}

\usepackage{xcolor}
\usepackage{graphicx}
\usepackage{booktabs}
\usepackage{tabularx}
\usepackage{longtable}
\usepackage{array}
\usepackage{makecell}
\usepackage{multirow}
\usepackage{enumitem}
\usepackage{fancyhdr}
\usepackage{titlesec}
\usepackage[most]{tcolorbox}
\usepackage{listings}
\usepackage{tikz}
\usetikzlibrary{arrows.meta,positioning,shapes.geometric,fit,backgrounds,calc}
\usepackage{caption}
\usepackage{amsmath,amssymb}
\usepackage{lastpage}
\usepackage{needspace}
\usepackage{fontspec}
\usepackage{hyperref}

% ============================================================
%  Fonts
% ============================================================
\setmainfont{Times New Roman}
\setsansfont{Segoe UI}
\setmonofont{Consolas}[Scale=0.92]
\setCJKmainfont{SimSun}[BoldFont=SimHei, ItalicFont=KaiTi]
\setCJKsansfont{Microsoft YaHei}
\setCJKmonofont{FangSong}[AutoFakeSlant=0.2]

% ============================================================
%  Colors
% ============================================================
\definecolor{primary}{HTML}{0B4F6C}
\definecolor{primarydark}{HTML}{073B4C}
\definecolor{accent}{HTML}{E76F51}
\definecolor{ink}{HTML}{1F2933}
\definecolor{muted}{HTML}{6B7280}
\definecolor{light}{HTML}{F4F7F9}
\definecolor{codebg}{HTML}{F7F8FA}
\definecolor{rulegray}{HTML}{D9DEE3}
\definecolor{success}{HTML}{2A9D8F}
\definecolor{warning}{HTML}{E9C46A}
\definecolor{danger}{HTML}{E63946}

% ============================================================
%  Hyperref
% ============================================================
\hypersetup{
  colorlinks=true,
  linkcolor=primary,
  urlcolor=accent,
  citecolor=primary,
  pdftitle={OpenCoxswain · 舵手 — 产品方案与 PRD},
  pdfauthor={OpenCoxswain 项目组},
  pdfsubject={本地优先的 Windows AI 总管},
  pdfkeywords={AI Agent, Orchestrator, Windows, Open Source, PRD},
  bookmarksnumbered=true,
  bookmarksopen=true,
  bookmarksopenlevel=1
}

% ============================================================
%  Page style
% ============================================================
\pagestyle{fancy}
\fancyhf{}
\fancyhead[L]{\small\sffamily\color{muted}\textcolor{accent}{\rule{0.55em}{0.55em}}\hspace{0.55em}OpenCoxswain · 舵手}
\fancyhead[R]{\small\sffamily\color{muted}产品方案与 PRD · v0.3}
\fancyfoot[C]{\small\sffamily\color{muted}\thepage\ /\ \pageref{LastPage}}
\renewcommand{\headrulewidth}{0.4pt}
\renewcommand{\headrule}{\hbox to\headwidth{\color{rulegray}\leaders\hrule height \headrulewidth\hfill}}

% ============================================================
%  Section styles
% ============================================================
\titleformat{\section}
  {\Large\bfseries\sffamily\color{primary}}
  {\thesection}{0.8em}{}
  [\vspace{2pt}{\color{rulegray}\titlerule[0.6pt]}]
\titleformat{\subsection}
  {\large\bfseries\sffamily\color{ink}}
  {\thesubsection}{0.7em}{}
\titleformat{\subsubsection}
  {\normalsize\bfseries\sffamily\color{ink}}
  {\thesubsubsection}{0.6em}{}

% Avoid orphaned headings at the bottom of a page
\let\ocsection\section
\renewcommand{\section}{\Needspace{8\baselineskip}\ocsection}
\let\ocsubsection\subsection
\renewcommand{\subsection}{\Needspace{5\baselineskip}\ocsubsection}
\let\ocsubsubsection\subsubsection
\renewcommand{\subsubsection}{\Needspace{4\baselineskip}\ocsubsubsection}

\setlength{\parindent}{2em}
\setlength{\parskip}{0.35em plus 0.1em minus 0.05em}
\linespread{1.32}
\setcounter{tocdepth}{2}

% Keep floats close to their references
\renewcommand{\topfraction}{0.9}
\renewcommand{\bottomfraction}{0.85}
\renewcommand{\textfraction}{0.06}
\renewcommand{\floatpagefraction}{0.8}
\setcounter{topnumber}{3}
\setcounter{bottomnumber}{3}
\setcounter{totalnumber}{5}

\setlist[itemize]{leftmargin=2em,itemsep=0.2em,topsep=0.3em}
\setlist[enumerate]{leftmargin=2.2em,itemsep=0.2em,topsep=0.3em}
\renewcommand{\arraystretch}{1.22}

\captionsetup{
  font=small,
  labelfont={bf,sf,color=primary},
  labelsep=quad,
  skip=6pt,
  hypcap=false
}

% ============================================================
%  Listings
% ============================================================
\lstset{
  basicstyle=\ttfamily\small,
  keywordstyle=\color{primary}\bfseries,
  stringstyle=\color{accent!85!black},
  commentstyle=\color{muted}\itshape,
  numberstyle=\tiny\color{muted},
  numbers=left,
  numbersep=8pt,
  frame=single,
  rulecolor=\color{rulegray},
  backgroundcolor=\color{codebg},
  breaklines=true,
  breakatwhitespace=false,
  columns=fullflexible,
  keepspaces=true,
  showstringspaces=false,
  tabsize=2,
  captionpos=b,
  xleftmargin=1.2em,
  framexleftmargin=1.2em,
  aboveskip=0.8em,
  belowskip=0.8em
}

% ============================================================
%  Helpers
% ============================================================
\newcommand{\code}[1]{\texttt{#1}}
\newcommand{\term}[1]{\textbf{#1}}
\newcommand{\pzero}{\textcolor{accent}{\textbf{P0}}}
\newcommand{\pone}{\textcolor{primary}{\textbf{P1}}}
\newcommand{\ptwo}{\textcolor{muted}{\textbf{P2}}}
\newcommand{\good}{\textcolor{success}{\textbf{可用}}}
\newcommand{\bad}{\textcolor{danger}{\textbf{冲突}}}
\newcommand{\warn}{\textcolor{warning!80!black}{\textbf{注意}}}

\newtcolorbox{keybox}[1][]{
  colback=light,
  colframe=primary,
  boxrule=0.4pt,
  arc=2pt,
  left=8pt,
  right=8pt,
  top=6pt,
  bottom=6pt,
  fonttitle=\sffamily\bfseries,
  coltitle=white,
  colbacktitle=primary,
  title={#1}
}

\newtcolorbox{notebox}[1][]{
  colback=white,
  colframe=rulegray,
  boxrule=0.4pt,
  arc=2pt,
  left=8pt,
  right=8pt,
  top=6pt,
  bottom=6pt,
  fonttitle=\sffamily\bfseries\color{ink},
  coltitle=ink,
  colbacktitle=light,
  title={#1}
}

% ============================================================
%  Document
% ============================================================
\begin{document}

% ------------------------------------------------------------
%  Cover
% ------------------------------------------------------------
\begin{titlepage}
\thispagestyle{empty}
\begin{tikzpicture}[remember picture, overlay]
  \fill[primary] (current page.north west) rectangle ([yshift=-2.8cm]current page.north east);
  \fill[accent] ([yshift=-2.8cm]current page.north west) rectangle ([yshift=-2.95cm]current page.north east);
  \node[anchor=west, text=white, font=\sffamily\bfseries\large]
    at ([xshift=2.3cm,yshift=-1.45cm]current page.north west)
    {PRODUCT REQUIREMENTS DOCUMENT};
  \node[anchor=west, text=white, font=\sffamily\small]
    at ([xshift=2.3cm,yshift=-2.15cm]current page.north west)
    {OpenCoxswain · 舵手};
  % small crew motif
  \begin{scope}[shift={([xshift=-4.0cm,yshift=-6.2cm]current page.north east)}]
    \node[circle, draw=primary, fill=light, minimum size=0.9cm, font=\scriptsize\sffamily] (c) at (0,0) {总管};
    \node[circle, draw=accent, fill=white, minimum size=0.65cm, font=\tiny\sffamily] (w1) at (2.2,0.9) {员工};
    \node[circle, draw=accent, fill=white, minimum size=0.65cm, font=\tiny\sffamily] (w2) at (2.2,0) {员工};
    \node[circle, draw=accent, fill=white, minimum size=0.65cm, font=\tiny\sffamily] (w3) at (2.2,-0.9) {员工};
    \draw[-{Latex[length=1.6mm]}, draw=muted] (c) -- (w1);
    \draw[-{Latex[length=1.6mm]}, draw=muted] (c) -- (w2);
    \draw[-{Latex[length=1.6mm]}, draw=muted] (c) -- (w3);
  \end{scope}
\end{tikzpicture}

\vspace*{3.6cm}

{\fontsize{46}{54}\selectfont\sffamily\bfseries\color{primary} OpenCoxswain\par}
\vspace{0.35cm}
{\fontsize{30}{36}\selectfont\sffamily\bfseries\color{ink} 舵手\par}
\vspace{0.55cm}
{\Large\sffamily\color{accent} The coxswain doesn't row.\par}
{\Large\sffamily\color{accent} 舵手不划桨。\par}

\vspace{1.0cm}
{\large\color{ink}
本地优先的 Windows AI 总管\\[0.25em]
你只对一个总管说目标，它带一队 AI 员工把活干完。\par}

\vspace{0.6cm}
\begin{minipage}{0.82\linewidth}
\small\color{muted}
这不是“再做一个更强的 Agent”，而是把项目经理 / Tech Lead 的职能产品化：
拆活、派活、盯进度、验收、记经验、向你汇报。
\end{minipage}

\vfill
\noindent\begin{tabularx}{\linewidth}{@{}l X@{}}
\toprule
\textbf{版本} & v0.3（讨论稿） \\
\textbf{日期} & 2026-11-09 \\
\textbf{状态} & 待评审，范围未冻结 \\
\textbf{首发平台} & Windows 10/11 x64 \\
\textbf{首发执行体} & OpenCode；Claude Code / Codex 后续接入 \\
\textbf{许可证} & Apache-2.0（建议，待最终确认） \\
\bottomrule
\end{tabularx}

\vspace{0.8cm}
{\footnotesize\sffamily\color{muted} OpenCoxswain 项目组 · 2026}
\end{titlepage}
\setcounter{page}{2}

% ------------------------------------------------------------
%  Document info
% ------------------------------------------------------------
\section*{文档信息}
\addcontentsline{toc}{section}{文档信息}

\noindent\begin{minipage}{\linewidth}
\centering
\small
\begin{tabularx}{\linewidth}{@{}p{3.2cm} X@{}}
\toprule
\textbf{产品名} & OpenCoxswain（中文名：舵手） \\
\textbf{一句话定位} & 本地优先的 Windows AI 总管：你只对一个总管说目标，它带一队 AI 员工把活干完。 \\
\textbf{文档版本} & v0.3（讨论稿） \\
\textbf{日期} & 2026-11-09 \\
\textbf{状态} & 待评审；v0 范围可讨论，未冻结 \\
\textbf{首发平台} & Windows 10/11 x64 \\
\textbf{技术形态} & Tauri 2 + React + TypeScript；Node/TS sidecar；SQLite + FTS5 \\
\textbf{首发执行体} & OpenCode（第一个 RuntimeAdapter） \\
\textbf{记忆方案} & Markdown 文件 + SQLite FTS5；v1 再考虑 embedding \\
\textbf{许可证} & Apache-2.0（建议）；AGPL-3.0 作为备选 \\
\textbf{商业模式} & 暂不考虑商业化；BYO API Key；本地优先 \\
\textbf{文档目的} & 对齐产品方向、范围、关键机制、v0 验收标准与开源策略 \\
\bottomrule
\end{tabularx}
\end{minipage}

\subsection*{修订记录}
\noindent\begin{minipage}{\linewidth}
\centering
\small
\begin{tabularx}{\linewidth}{@{}p{1.5cm} p{2.4cm} X@{}}
\toprule
\textbf{版本} & \textbf{日期} & \textbf{变更} \\
\midrule
v0.1 & 2026-11-09 & 初版 PRD：定位、四层架构、v0/v1/v2、关键机制。 \\
v0.2 & 2026-11-09 & 确认 Windows 优先、OpenCode 首发、Markdown + FTS5、本地优先、不做商业化。 \\
v0.3 & 2026-11-09 & 产品定名 OpenCoxswain；补充 OpenSuperAgent 命名冲突分析、开源策略、Windows 专项验收。 \\
\bottomrule
\end{tabularx}
\end{minipage}

\subsection*{阅读指南}
\begin{itemize}
  \item 只想快速判断方向：读第 \ref{sec:summary} 节（摘要）与第 \ref{sec:roadmap} 节（版本规划与验收）。
  \item 关心产品范围：读第 \ref{sec:users} 至第 \ref{sec:features} 节。
  \item 关心技术落地：读第 \ref{sec:mechanisms} 与第 \ref{sec:architecture} 节。
  \item 关心开源与命名：读第 \ref{sec:opensource} 与第 \ref{sec:naming} 节。
\end{itemize}

\newpage
\tableofcontents
\newpage

% ============================================================
%  1. 摘要
% ============================================================
\section{摘要}
\label{sec:summary}

OpenCoxswain（中文名：舵手）是一个\textbf{本地优先的 Windows AI 总管}。用户只对一个常驻的“总管”说话；总管不直接动手，而是把目标拆成结构化任务单，派给多个独立的 AI 员工，盯进度、做验收，最后把\textbf{待决策项、未通过的检查、交付产物}汇总回来。

它的核心不是“再做一个更强的 Agent”，而是把项目经理 / Tech Lead 的职能产品化：\textbf{规划与执行解耦}。执行能力已经商品化，真正稀缺的是“谁来拆、谁来盯、谁来记、谁来验、谁来控成本”。

\begin{keybox}[一句话记住这个产品]
你只跟一个“人”说话。这个“人”不干活，它管理一队干活的“人”。
\end{keybox}

\subsection*{v0 目标}
v0 只做一个小闭环：

\begin{itemize}
  \item 单项目、一个总管、两个员工（实施 + 校验）。
  \item 结构化任务单、执行清单确认、DAG 拆解（2--4 个子任务）。
  \item \code{git worktree} 隔离、状态机、极简看板。
  \item 预算上限、危险动作人审、事件日志、可回放。
  \item 汇总固定三块：待决策 / 未通过 / 产物。
\end{itemize}

\textbf{验收标准}：一句真实需求进去，两张任务卡并行跑起来，最后拿到一份带待决策项的汇报；全程可回放，总花费有上限，无人工接管代码。

% ============================================================
%  2. 背景与机会
% ============================================================
\section{背景与机会}

\subsection{执行能力正在商品化}

Claude Code、Codex、Trae SOLO、OpenHands 等工具已经能写代码、跑测试、开 PR、做调研。执行层的竞争非常激烈，能力差距在缩小。真正没有被产品化的是\textbf{管理}：

\noindent\begin{minipage}{\linewidth}
\centering
\captionof{table}{瓶颈已经从“能不能做”转向“谁来管理”}
\small
\begin{tabularx}{\linewidth}{@{}p{4.2cm} X@{}}
\toprule
\textbf{瓶颈} & \textbf{现状} \\
\midrule
谁来拆任务 & 用户自己拆，或者 Agent 边做边猜。 \\
谁来盯多个会话 & 用户自己开一堆窗口来回切。 \\
谁来记上下文 & 每个会话重新讲一遍背景。 \\
谁来验收 & 用户自己看 diff、跑测试、判断能不能交。 \\
谁来控制成本 & 用户自己盯着 token 和并发。 \\
\bottomrule
\end{tabularx}
\end{minipage}

\subsection{现有工具的缺口}

现有工具的并行能力仍是\textbf{会话内}的：Subagent、Agent Teams 都要用户自己当调度器。用户承担了项目经理的工作，却没有项目经理的视图、记忆和工具。

\begin{itemize}
  \item \textbf{没有管理视图}：不知道谁在干什么、卡在哪、花了多少。
  \item \textbf{没有组织记忆}：换项目、换天就失忆，重复讲背景。
  \item \textbf{没有硬约束}：总管很容易“顺手自己干”，多员工退化成单 Agent。
  \item \textbf{没有验收闭环}：员工说“完成了”，用户还得自己重跑一遍。
\end{itemize}

\subsection{机会}

把“管理”从会话里拿出来，做成一个\textbf{常驻的人 + 一套组织记忆 + 一个管理工作台}。执行体可以插拔，用户自带 API Key，代码和记忆默认不出本机。

产品策略：

\begin{enumerate}
  \item \textbf{切入点}：单项目、编码场景、一个总管 + 两个员工，先跑通闭环。
  \item \textbf{护城河}：记忆质量 + 工作台使用习惯 + 执行体中立。
  \item \textbf{扩张路径}：编码 $\rightarrow$ 调研/写作/数据处理 $\rightarrow$ 多项目 $\rightarrow$ 小团队协作。
\end{enumerate}

% ============================================================
%  3. 目标用户
% ============================================================
\section{目标用户}
\label{sec:users}

\subsection{主用户：独立开发者 / 小团队 Tech Lead}

\begin{itemize}
  \item 已经在用 Claude Code、Cursor、Trae、Codex 等 AI 编码工具。
  \item 同时推进 2--5 个项目，经常“自己就是那个调度器”。
  \item 痛点：上下文切换、重复讲背景、不知道子任务卡在哪、并行烧钱失控。
  \item 愿意为“少管一点”付费的不是钱，而是注意力和时间。
\end{itemize}

\subsection{次用户：AI 重度创作者}

需要“调研 $\rightarrow$ 成文 $\rightarrow$ 校验”这类非纯编码流程。v1 再覆盖，v0 不分散精力。

\subsection{v0 明确不服务}

\begin{itemize}
  \item 完全没用过 AI 编码工具的新手。
  \item 大型企业团队、需要合规审批链的组织。
  \item 非技术用户的一站式自动化。
\end{itemize}

% ============================================================
%  4. 竞品与差异化
% ============================================================
\section{竞品与差异化}

\noindent\begin{minipage}{\linewidth}
\centering
\captionof{table}{竞品定位与我们的差异}
\small
\begin{tabularx}{\linewidth}{@{}p{2.6cm} p{2.8cm} X X@{}}
\toprule
\textbf{产品} & \textbf{形态} & \textbf{强项} & \textbf{缺口 / 我们的差异} \\
\midrule
TRAE 宝 & 字节 Trae 的个人助理 & 常驻人格、跨项目记忆、工作台体感 & 绑定自家生态；我们做本地优先、执行体可插拔、记忆可读可改。 \\
Trae SOLO / Claude Code Subagent / Codex Agent Teams & 会话内多 Agent & 执行强、生态成熟 & 仍是用户当调度器；我们做会话之上的管理层。 \\
Devin & 云端自主工程师 & 端到端交付 & 云端、成本高、代码出本机；我们做本地/自带 Key/成本可控。 \\
CrewAI / LangGraph & 开发框架 & 编排灵活 & 面向开发者，不是终端产品；我们做开箱即用的产品体验。 \\
Cursor & AI IDE & 编辑体验 & 单会话为中心；我们不替换编辑器，做上面的管家。 \\
\bottomrule
\end{tabularx}
\end{minipage}

\begin{keybox}[定位语]
本地优先、执行体可插拔、记忆可读可改的 AI 项目管家。
\end{keybox}

% ============================================================
%  5. 产品原则
% ============================================================
\section{产品原则（不可妥协）}

\begin{enumerate}
  \item \textbf{单入口}：用户只跟总管说话，不直接指挥员工。
  \item \textbf{总管不动手}：系统层清零它的文件写与 shell 权限，只留元操作。
  \item \textbf{状态在代码里}：任务树、状态机、预算、审批全部持久化，不靠模型记忆。
  \item \textbf{上下文隔离}：员工只拿自己那片任务，日志留在沙箱，只回摘要和产物。
  \item \textbf{一切可回放}：每次派单、每次工具调用、每次验收都有事件记录。
  \item \textbf{随时可接管}：用户可以暂停、终止、改任务单、直接看员工会话。
  \item \textbf{简单任务走轻路径}：不为“多智能体”而多智能体，开会话比干活快时就直接干。
  \item \textbf{本地优先}：代码、密钥、记忆默认不出本机；云端执行是可选能力。
\end{enumerate}

\begin{keybox}[硬约束：总管不动手]
总管的工具注册表里没有文件读写、shell、网络、git。它只能调用：
\code{spawn} / \code{assign} / \code{query} / \code{collect} / \code{ask\_human} / \code{read\_memory} / \code{write\_memory}。
这些能力由调度器校验状态、权限和预算后再执行。靠提示词约束它一定会顺手自己干。
\end{keybox}

% ============================================================
%  6. 核心概念
% ============================================================
\section{核心概念}

\noindent\begin{minipage}{\linewidth}
\centering
\captionof{table}{核心概念与职责边界}
\small
\begin{tabularx}{\linewidth}{@{}p{3.0cm} X@{}}
\toprule
\textbf{概念} & \textbf{定义} \\
\midrule
总管（Steward） & 常驻人格，负责澄清、拆解、派发、追问、验收、汇总；不直接执行。 \\
员工（Worker） & 一次任务实例：独立会话 + 独立上下文 + 独立 worktree + 绑定角色/工具/模型。 \\
角色模板（Role） & 调研员、实施员、校验员、文员等，定义工具权限与输出格式。 \\
任务单（Task Contract） & 结构化 JSON：目标、输入、交付物、约束、验收标准、依赖、预算。 \\
任务树（Task DAG） & 任务之间的依赖与并行关系。 \\
黑板（Blackboard） & 共享产物区：文件、diff、报告、链接；员工不自由聊天，只通过产物交接。 \\
工作台（Workbench） & 对话 + 任务树 + 员工看板 + 产物 + 待决策 + 花费。 \\
记忆（Memory） & 全局 \code{user\_profile.md} + 项目 \code{project\_memory.md} + 决策/教训。 \\
人审点（Approval Gate） & 删除、推送、发布、付款、外发等不可逆动作前的强制确认。 \\
\bottomrule
\end{tabularx}
\end{minipage}

% ============================================================
%  7. 用户旅程
% ============================================================
\section{用户旅程}

\subsection{首次使用}

\begin{enumerate}
  \item 选择本地 git 仓库。
  \item 给总管起名、选头像、选语气（简洁/详细）和主动性（先问再做/先做后报）。
  \item 设置预算：单任务上限、每日上限、最大并发。
  \item 总管扫描项目结构，生成 \code{project\_memory.md} 草稿，用户确认。
  \item 跑一个“只读调研”小任务，验证 runtime 和权限。
\end{enumerate}

\subsection{日常派活}

\begin{notebox}[一个典型请求]
用户：「把设置页夜间模式切换后不生效的问题修了。」
\end{notebox}

\begin{enumerate}
  \item 总管最多问 1--3 个阻塞性问题；不阻塞就带假设继续。
  \item 产出\textbf{任务单}：目标、涉及文件、步骤、验收标准、预算、风险点。
  \item 用户确认“开工”或“改一下”。
  \item 总管拆成 2--4 个子任务，形成 DAG；独立的并行，有依赖的串行。
  \item 员工在各自 worktree 执行；工作台实时显示状态、当前动作、耗时、花费。
  \item 校验员独立复跑测试、检查 diff。
  \item 总管汇总三块：\textbf{待决策项 / 未通过的检查 / 交付产物}。
  \item 用户点开 diff，选择通过或打回；打回则生成返工任务。
\end{enumerate}

\subsection{中途干预}

\begin{itemize}
  \item 暂停某个员工，保留现场。
  \item 改任务单（加约束、换验收标准）。
  \item 加人/减人，调整并发。
  \item 直接打开员工会话接管。
  \item 降低预算或直接熔断。
\end{itemize}

\subsection{复盘与记忆}

\begin{itemize}
  \item 任务验收后，记忆管理员提取 0--3 条候选记忆。
  \item 用户确认后写入 \code{project\_memory.md} / \code{decisions.md} / \code{lessons.md}。
  \item 下次任务开始时，总管展示“我参考了这些记忆”，可点开纠正。
\end{itemize}

\subsection{跨项目（v2）}

\begin{itemize}
  \item 全局助理能看各项目摘要与历史决策。
  \item 对每条可复用经验标注：\textbf{可复用 / 仅限原项目 / 需要重新确认}。
  \item 典型提示：“上次类似方案是 X，但约束 Y 只适用于旧项目。”
\end{itemize}

% ============================================================
%  8. 功能需求
% ============================================================
\section{功能需求}
\label{sec:features}

\subsection{人格层}

\noindent\begin{minipage}{\linewidth}
\centering
\captionof{table}{人格层功能}
\small
\begin{tabularx}{\linewidth}{@{}p{2.6cm} X p{1.6cm}@{}}
\toprule
\textbf{功能} & \textbf{说明与验收} & \textbf{优先级} \\
\midrule
工牌 & 名字、头像、一句话人设；可改名、可换头像。 & \pzero \\
单入口对话 & 所有指令只对总管说；界面上没有“直接指挥员工”的入口。 & \pzero \\
语气配置 & 简洁/详细、主动/被动；配置生效，不改变硬约束。 & \pzero \\
状态播报 & “我在等小李的测试结果”；任务运行中至少每 30 秒有一次状态更新。 & \pzero \\
汇报格式 & 固定三块：待决策 / 未通过 / 产物；每次任务结束都按此格式。 & \pzero \\
托盘/快捷键 & 随时呼出；全局快捷键可配置。 & \pone \\
桌宠 & 可选形象；不影响主流程。 & \ptwo \\
IM 入口 & 飞书/Telegram/Discord bot；能收任务、能审批、能看状态。 & \ptwo \\
\bottomrule
\end{tabularx}
\end{minipage}

\subsection{任务单与编排}

\noindent\begin{minipage}{\linewidth}
\centering
\captionof{table}{任务单与编排功能}
\small
\begin{tabularx}{\linewidth}{@{}p{2.6cm} X p{1.6cm}@{}}
\toprule
\textbf{功能} & \textbf{说明与验收} & \textbf{优先级} \\
\midrule
目标澄清 & 最多 3 个阻塞性问题，其余带假设；80\% 的日常任务 1 分钟内出任务单。 & \pzero \\
任务单生成 & 结构化字段完整；缺字段不允许进入执行。 & \pzero \\
执行清单确认 & 文件/步骤/验收/预算/风险；未确认不能 spawn 员工。 & \pzero \\
自动拆解 & 2--5 个子任务 + DAG；拆解结果可视化、可手改。 & \pzero \\
派单与依赖 & 独立并行、依赖串行；有写冲突的任务不并行。 & \pzero \\
进度追踪 & 状态机驱动；UI 延迟小于 1 秒。 & \pzero \\
验收与打回 & 按验收标准逐条判定；每条标准有通过/失败/未验证。 & \pzero \\
汇总 & 三块固定格式；不允许只回“完成了”。 & \pzero \\
任务模板 & 调研→成文→校验；修 bug→测试→review；模板可另存、可复用。 & \pone \\
失败重试/升级 & 原员工重试→换强模型→拆任务→问人；每级有次数上限。 & \pone \\
定时/触发任务 & cron、文件变更、webhook；可暂停、可手动触发。 & \ptwo \\
\bottomrule
\end{tabularx}
\end{minipage}

\subsection{员工执行层}

\noindent\begin{minipage}{\linewidth}
\centering
\captionof{table}{员工执行层功能}
\small
\begin{tabularx}{\linewidth}{@{}p{2.6cm} X p{1.6cm}@{}}
\toprule
\textbf{功能} & \textbf{说明与验收} & \textbf{优先级} \\
\midrule
角色模板 & 调研员/实施员/校验员/文员；每个角色有工具白名单。 & \pzero \\
独立 worktree & 每任务一个分支/目录；员工之间文件不互相污染。 & \pzero \\
独立会话 & 独立上下文窗口；员工看不到彼此日志。 & \pzero \\
最小权限 & 按角色绑定工具；校验员默认只读。 & \pzero \\
摘要回传 & 只回摘要 + 产物链接；原始日志不进总管上下文。 & \pzero \\
日志留存 & 本地可查、可导出；任务结束仍可回放。 & \pzero \\
并发上限 & 全局 + 项目级；超限排队，不静默丢弃。 & \pzero \\
暂停/终止/接管 & 保留现场；暂停后可从断点继续。 & \pzero \\
合并与冲突 & 顺序合并，冲突转人审；冲突不自动强推。 & \pzero \\
容器沙箱 & 可选 Docker 隔离；默认关闭，开启后工具受限。 & \ptwo \\
\bottomrule
\end{tabularx}
\end{minipage}

\subsection{记忆层}

\noindent\begin{minipage}{\linewidth}
\centering
\captionof{table}{记忆层功能}
\small
\begin{tabularx}{\linewidth}{@{}p{2.6cm} X p{1.6cm}@{}}
\toprule
\textbf{功能} & \textbf{说明与验收} & \textbf{优先级} \\
\midrule
全局记忆 & \code{user\_profile.md}；跨项目可用，用户可编辑。 & \pone \\
项目记忆 & \code{project\_memory.md}；只在该项目加载。 & \pone \\
决策日志 & \code{decisions.md}；记录“为什么这么定”。 & \pone \\
教训库 & \code{lessons.md}；记录踩坑与规避方式。 & \pone \\
写入门槛 & 任务验收后才提取；敏感信息、一次性指令不写；用户可拒绝候选记忆。 & \pone \\
混合检索 & 关键词 + 向量 + 时间衰减；任务开始展示命中的记忆。 & \pone \\
冲突处理 & 新覆盖旧、项目覆盖全局、冲突上报；不静默覆盖。 & \pone \\
复用判断 & 可复用/项目绑定/需重问；跨项目时强制标注。 & \pone \\
记忆检查器 & 查看、编辑、删除、标注过期；一条记忆可追溯来源任务。 & \pone \\
隐私 & 密钥、账号、个人信息默认不写；写入前脱敏扫描。 & \pone \\
\bottomrule
\end{tabularx}
\end{minipage}

\subsection{工作台}

\noindent\begin{minipage}{\linewidth}
\centering
\begin{tikzpicture}[
  panel/.style={draw=primary, fill=light, rounded corners=2pt, minimum height=5.2cm, align=left, font=\footnotesize\sffamily, inner sep=8pt},
  label/.style={font=\small\sffamily\bfseries\color{primary}}
]
\node[panel, minimum width=4.0cm] (left) at (0,0) {\textbf{项目 / 任务}\\[4pt]$\blacktriangleright$ 当前项目\\[2pt]$\quad\blacktriangleright$ 任务树\\[2pt]$\quad\blacktriangleright$ 历史任务\\[2pt]$\quad\blacktriangleright$ 记忆};
\node[panel, minimum width=6.2cm] (mid) at (4.5,0) {\textbf{总管对话}\\[4pt]你：修一下夜间模式\\[2pt]总管：任务单如下……\\[4pt]\fbox{确认开工}\quad\fbox{改一下}};
\node[panel, minimum width=5.0cm] (right) at (10.0,0) {\textbf{员工看板}\\[4pt]\fbox{\begin{minipage}{4.2cm}\footnotesize 实施员 · 运行中\\当前：改 Theme.kt\\耗时 03:12 / \$0.42\end{minipage}}\\[4pt]\fbox{\begin{minipage}{4.2cm}\footnotesize 校验员 · 等待依赖\end{minipage}}};
\node[draw=rulegray, fill=white, minimum width=15.4cm, minimum height=1.0cm, font=\footnotesize\sffamily, anchor=north] at (5.0,-3.0) {标签：任务树 \quad 产物 \quad 待决策 \quad 记忆 \quad 花费 \quad 时间线};
\end{tikzpicture}
\captionof{figure}{工作台主界面信息架构示意}
\label{fig:workbench}
\end{minipage}

\noindent\begin{minipage}{\linewidth}
\centering
\captionof{table}{工作台必备视图}
\small
\begin{tabularx}{\linewidth}{@{}p{2.8cm} X p{1.6cm}@{}}
\toprule
\textbf{视图} & \textbf{说明} & \textbf{优先级} \\
\midrule
任务树/DAG & 依赖关系、关键路径、阻塞点。 & \pzero \\
员工卡 & 角色、状态、当前动作、进度、耗时、花费、暂停/终止/查看会话。 & \pzero \\
产物区 & diff、文件、报告、链接；支持逐文件审阅。 & \pzero \\
待决策收件箱 & 需要用户拍板的事项，按紧急度排序。 & \pone \\
花费 & 按任务/员工/模型拆分。 & \pone \\
时间线/回放 & 每个事件可展开，支持导出。 & \pzero \\
搜索 & 按任务、文件、记忆、产物搜索。 & \pone \\
\bottomrule
\end{tabularx}
\end{minipage}

\subsection{安全、权限与预算}

\noindent\begin{minipage}{\linewidth}
\centering
\captionof{table}{安全、权限与预算}
\small
\begin{tabularx}{\linewidth}{@{}p{3.0cm} X@{}}
\toprule
\textbf{能力} & \textbf{要求} \\
\midrule
总管工具白名单 & 只注册 \code{spawn / assign / query / collect / ask\_human / read\_memory / write\_memory}。 \\
员工最小权限 & 角色模板决定工具；校验员默认只读。 \\
人审门 & 删除、\code{git push}、发布、付款、对外发送、安装系统级依赖。 \\
预算熔断 & 单任务 70\% 警告、100\% 硬停；项目日预算同理。 \\
并发上限 & 默认 2，可配置；超限排队。 \\
审计日志 & 谁在什么时候做了什么、花了多少，append-only。 \\
密钥脱敏 & 不进入 prompt；日志输出前扫描。 \\
注入防护 & 员工输出视为不可信；总管只读结构化摘要；仓库内容不得改变系统规则。 \\
\bottomrule
\end{tabularx}
\end{minipage}

% ============================================================
%  9. 关键机制
% ============================================================
\section{关键机制}
\label{sec:mechanisms}

\subsection{任务状态机}

\noindent\begin{minipage}{\linewidth}
\centering
\begin{tikzpicture}[
  state/.style={draw=primary, fill=light, rounded corners=2pt, minimum width=1.72cm, minimum height=0.72cm, font=\footnotesize\sffamily, align=center},
  terminal/.style={state, fill=primary!12},
  arrow/.style={-{Latex[length=1.8mm]}, draw=muted, thick},
  back/.style={-{Latex[length=1.6mm]}, draw=muted, thick, dashed},
  bidir/.style={{Latex[length=1.5mm]}-{Latex[length=1.5mm]}, draw=muted, thick},
  lbl/.style={font=\tiny\sffamily, text=muted, inner sep=1.6pt}
]
\node[terminal] (draft) {draft};
\node[state, right=0.36cm of draft] (planning) {planning};
\node[state, right=0.36cm of planning] (approval) {awaiting\\approval};
\node[state, right=0.36cm of approval] (queued) {queued};
\node[state, right=0.36cm of queued] (running) {running};
\node[state, right=0.36cm of running] (review) {awaiting\\review};
\node[terminal, right=0.36cm of review] (done) {done};

\draw[arrow] (draft) -- (planning);
\draw[arrow] (planning) -- (approval);
\draw[arrow] (approval) -- (queued);
\draw[arrow] (queued) -- (running);
\draw[arrow] (running) -- (review);
\draw[arrow] (review) -- (done);

\draw[back] (approval) to[out=120,in=60] node[lbl,above] {打回} (planning);
\draw[back] (review) to[out=120,in=60] node[lbl,above] {返工} (running);

\node[state, below=1.15cm of planning] (cancelled) {cancelled};
\node[state, below=1.15cm of queued] (paused) {paused};
\node[state, below=1.15cm of running] (blocked) {blocked};
\node[state, below=1.15cm of review] (failed) {failed};

\draw[back] (planning.south) -- node[lbl,right] {取消} (cancelled.north);
\draw[bidir] (queued.south) -- node[lbl,right] {暂停 / 继续} (paused.north);
\draw[bidir] (running.south) -- node[lbl,right] {阻塞 / 恢复} (blocked.north);
\draw[back] (running.south east) to[out=0,in=90] node[lbl,right] {失败} (failed.north);
\draw[back] (failed.south) -- ++(0,-0.55) -- node[lbl,below] {重试} ([yshift=-0.55cm]queued.south) -- (queued.south);
\end{tikzpicture}
\captionof{figure}{任务状态机（简化）}
\label{fig:state-machine}
\end{minipage}

\noindent\begin{minipage}{\linewidth}
\centering
\captionof{table}{状态机规则}
\small
\begin{tabularx}{\linewidth}{@{}p{3.2cm} p{2.8cm} X@{}}
\toprule
\textbf{规则} & \textbf{触发者} & \textbf{说明} \\
\midrule
进入执行 & 用户 & 只有用户能触发 \code{awaiting\_approval} $\rightarrow$ \code{queued}（v0）；v1 可授权“预算内自动开工”。 \\
状态推进 & 调度器 & 员工不能改自己的状态，只能上报事件；状态由调度器推进。 \\
阻塞 & 员工/系统 & \code{blocked} 必须带原因：依赖失败、权限不足、冲突、预算、需要人审。 \\
完成 & 调度器 & \code{done} 必须所有验收标准都有结论。 \\
暂停/取消 & 用户 & 任意状态可暂停；暂停后回到上一状态；取消不可恢复。 \\
\bottomrule
\end{tabularx}
\end{minipage}

\subsection{任务单契约}

\begin{lstlisting}[language=Java,caption={任务单 JSON 示例},label={lst:task-contract}]
{
  "task_id": "tsk_01",
  "parent_id": null,
  "project_id": "proj_coursemeow",
  "title": "修复夜间模式切换不生效",
  "goal": "设置页切换夜间模式后，主题立即生效并在重启后保持",
  "context_refs": [
    "memory://project/theme",
    "file://app/src/main/java/.../Theme.kt"
  ],
  "inputs": {
    "repro": "设置页 -> 外观 -> 夜间模式 -> 返回首页仍为亮色"
  },
  "deliverables": [
    { "type": "diff", "scope": "app/src/main/java/.../ui/theme/**" },
    { "type": "test_report", "format": "markdown" }
  ],
  "constraints": [
    "不改数据库结构",
    "不引入新依赖",
    "保持现有配色 token 命名"
  ],
  "acceptance_criteria": [
    { "id": "ac1", "type": "command", "desc": "单测通过",
      "check": "gradlew testDebugUnitTest" },
    { "id": "ac2", "type": "rubric",
      "desc": "切换后首页、设置页、小组件均使用夜间配色" },
    { "id": "ac3", "type": "human", "desc": "用户确认视觉无回归" }
  ],
  "dependencies": [],
  "write_scope": [
    "app/src/main/java/.../ui/theme/**"
  ],
  "assignee": {
    "role": "implementer",
    "runtime": "opencode",
    "model": "strong"
  },
  "budget": { "tokens": 200000, "usd": 2.0, "minutes": 30 },
  "review": { "required": true, "reviewer_role": "verifier" },
  "deadline": "2026-11-10T18:00:00+08:00"
}
\end{lstlisting}

\noindent\begin{minipage}{\linewidth}
\centering
\captionof{table}{验收标准的四种类型}
\small
\begin{tabularx}{\linewidth}{@{}p{2.4cm} p{3.6cm} X@{}}
\toprule
\textbf{类型} & \textbf{谁来判} & \textbf{例子} \\
\midrule
\code{command} & 系统 & 跑测试、lint、构建。 \\
\code{diff} & 系统/校验员 & 只允许改指定目录、不得提交密钥。 \\
\code{rubric} & 校验员（LLM + 评分表） & 视觉一致性、文案规范。 \\
\code{human} & 用户 & 最终体验、产品取舍。 \\
\bottomrule
\end{tabularx}
\end{minipage}

\subsection{总管工具白名单}

\noindent\begin{minipage}{\linewidth}
\centering
\captionof{table}{总管可调用的元操作}
\small
\begin{tabularx}{\linewidth}{@{}p{3.0cm} p{4.6cm} p{3.6cm} X@{}}
\toprule
\textbf{工具} & \textbf{输入} & \textbf{输出} & \textbf{说明} \\
\midrule
\code{spawn\_worker} & role, runtime, model, task\_id & worker\_id & 创建员工实例。 \\
\code{assign\_task} & task\_id, worker\_id & ack & 派单，校验依赖与写域。 \\
\code{query\_worker} & worker\_id & status\_summary & 查状态，不读原始日志。 \\
\code{collect\_result} & worker\_id & summary + artifacts & 收结果，走结构化 schema。 \\
\code{ask\_human} & question, options, blocking & answer & 需要拍板时阻塞等待。 \\
\code{read\_memory} & scope, query & entries & 读记忆，带来源。 \\
\code{write\_memory} & entry, scope & ack & 写候选记忆，需通过门槛。 \\
\bottomrule
\end{tabularx}
\end{minipage}

\begin{notebox}[没有的能力]
文件读写、shell、网络、git、安装依赖、直接调用模型。这些都在员工侧，由角色模板和最小权限控制。
\end{notebox}

\subsection{记忆读写规则}

\subsubsection*{写入门槛（全部满足才写）}

\begin{enumerate}
  \item 任务已验收通过。
  \item 内容可复用，不是一次性指令。
  \item 不含密钥、账号、个人敏感信息。
  \item 有明确来源任务。
  \item 不与现有记忆重复；冲突时走覆盖/并存策略。
\end{enumerate}

\subsubsection*{读取顺序}

\begin{enumerate}
  \item 全局 \code{user\_profile.md}（永远加载）。
  \item 项目 \code{project\_memory.md} 摘要。
  \item 与当前任务相关的 top 5--10 条（关键词 + 向量 + 时间衰减）。
  \item 展示“我参考了这些记忆”，用户可点开纠正。
\end{enumerate}

\noindent\begin{minipage}{\linewidth}
\centering
\captionof{table}{跨项目复用的三种判断}
\small
\begin{tabularx}{\linewidth}{@{}p{3.2cm} X@{}}
\toprule
\textbf{判断} & \textbf{处理} \\
\midrule
可复用 & 直接引用，标注来源。 \\
仅限原项目 & 不引用，只提示“旧项目有类似方案”。 \\
需要重新确认 & 引用但必须问用户。 \\
\bottomrule
\end{tabularx}
\end{minipage}

\subsection{汇总模板}

\begin{lstlisting}[caption={任务结束时的固定汇总格式},label={lst:summary}]
## 待你决策
- [ ] 是否接受方案 A（影响：...）
- [ ] 是否允许推送分支（涉及 3 个 commit）

## 未通过的检查
- [x] ac2：小组件夜间配色未覆盖（证据：截图/日志链接）
- [!] ac3：需要你人工确认视觉

## 交付产物
- 分支：feat/night-mode-fix
- diff：3 个文件，+48 / -12
- 测试报告：artifacts/tsk_01/test-report.md
- 花费：$0.87 / 预算 $2.00
\end{lstlisting}

\subsection{失败处理矩阵}

\noindent\begin{minipage}{\linewidth}
\centering
\captionof{table}{失败处理矩阵}
\small
\begin{tabularx}{\linewidth}{@{}p{2.8cm} X X X@{}}
\toprule
\textbf{失败类型} & \textbf{第一处理} & \textbf{第二处理} & \textbf{最后处理} \\
\midrule
员工报错 & 原员工带错误上下文重试 1 次 & 换更强模型 & 拆小任务 / 问人。 \\
超时 & 2 倍预估时间暂停并汇报 & 降范围 & 问人。 \\
测试不过 & 员工自修 1 次 & 校验员给修复建议 & 问人。 \\
合并冲突 & 标记 blocked & 启动集成员工 & 问人。 \\
预算超限 & 70\% 警告 & 100\% 硬停 & 问人追加。 \\
权限不足 & 请求人审 & 换只读方案 & 取消任务。 \\
\bottomrule
\end{tabularx}
\end{minipage}

% ============================================================
%  10. 技术架构
% ============================================================
\section{技术架构}
\label{sec:architecture}

\subsection{总体形态}

\noindent\begin{minipage}{\linewidth}
\centering
\begin{tikzpicture}[
  box/.style={draw=primary, fill=light, rounded corners=3pt, minimum width=8.0cm, minimum height=0.8cm, align=center, font=\small\sffamily},
  smallbox/.style={draw=rulegray, fill=white, rounded corners=2pt, minimum width=3.4cm, minimum height=0.72cm, align=center, font=\footnotesize\sffamily},
  arrow/.style={-{Latex[length=2mm]}, draw=muted, thick}
]
\node[box] (ui) {桌面工作台（Tauri 2 + React + TypeScript）};
\node[box, below=0.7cm of ui] (daemon) {本地 Daemon（Node/TS sidecar）};
\node[smallbox, below=0.5cm of daemon, xshift=-4.3cm] (orch) {总管服务};
\node[smallbox, below=0.5cm of daemon] (sched) {调度器 + 状态机};
\node[smallbox, below=0.5cm of daemon, xshift=4.3cm] (mem) {记忆服务};
\node[draw=rulegray, dashed, rounded corners=3pt, fit=(orch)(sched)(mem), inner sep=7pt] (layer) {};
\node[font=\footnotesize\sffamily\color{muted}, anchor=south west] at (layer.north west) {编排层};
\node[box, below=1.0cm of layer] (runtime) {Runtime Adapters（OpenCode / Claude Code / Codex）};
\node[box, below=0.55cm of runtime] (wt) {git worktree + 独立会话 + 独立上下文};
\draw[arrow] (ui) -- node[right,font=\footnotesize\sffamily,text=muted] {WebSocket / SSE} (daemon);
\draw[arrow] (daemon) -- (runtime);
\draw[arrow] (runtime) -- (wt);
\end{tikzpicture}
\captionof{figure}{OpenCoxswain 总体架构}
\label{fig:architecture}
\end{minipage}

\textbf{关键决策}：UI 与执行解耦。Daemon 是产品核心，桌面 App、CLI、IM bot 都只是入口。这样 v2 加 IM/云端时不用重写。

\subsection{Windows 桌面技术栈}

\noindent\begin{minipage}{\linewidth}
\centering
\captionof{table}{Windows 桌面技术选型}
\small
\begin{tabularx}{\linewidth}{@{}p{3.0cm} p{4.0cm} X@{}}
\toprule
\textbf{层} & \textbf{选型} & \textbf{说明} \\
\midrule
桌面壳 & Tauri 2 & 安装包小、内存低、Windows 支持成熟；后续可扩 macOS/Linux。 \\
UI & React 18 + TypeScript & 生态成熟；工作台组件多，Web 技术迭代快。 \\
本地 Daemon & Node 22 + TypeScript & 与 OpenCode/Agent SDK 生态一致；作为 sidecar 随 App 启动。 \\
数据库 & SQLite + FTS5 & 单机零运维；任务、事件、记忆索引都放这里。 \\
运行时 & WebView2 & Windows 10/11 自带；安装包做检测与引导。 \\
打包 & NSIS / MSI & 默认 per-user 安装，避免管理员权限。 \\
自动更新 & Tauri Updater + GitHub Releases & 开源项目最顺；更新包签名。 \\
密钥存储 & Windows Credential Manager / DPAPI & API Key 不进明文配置，不进日志。 \\
\bottomrule
\end{tabularx}
\end{minipage}

\subsection{RuntimeAdapter 接口}

\begin{lstlisting}[language=Java,caption={RuntimeAdapter 接口（TypeScript 伪代码）},label={lst:adapter}]
interface RuntimeAdapter {
  readonly name: string;                 // "opencode" | "claude-code" | "codex"

  createSession(opts: SessionOptions): Promise<RuntimeSession>;

  send(
    session: RuntimeSession,
    task: TaskContract
  ): AsyncIterable<RuntimeEvent>;

  cancel(session: RuntimeSession): Promise<void>;

  collect(session: RuntimeSession): Promise<WorkerResult>;

  dispose(session: RuntimeSession): Promise<void>;
}

interface SessionOptions {
  taskId: string;
  worktree: string;          // 独立目录
  role: RoleTemplate;        // 工具白名单
  model: ModelConfig;        // provider / model / budget
  approvalPolicy: ApprovalPolicy;
}

interface WorkerResult {
  summary: string;           // 只回摘要
  artifacts: ArtifactRef[];  // 产物链接
  acceptance: AcceptanceResult[];
  cost: CostLedger;
  logsRef: string;           // 原始日志留在本地
}
\end{lstlisting}

\textbf{首发实现}：OpenCode。先用一个 spike 验证 Windows 下 \code{opencode serve} / SDK 的稳定性；如果不稳，退回 CLI 包装，或把 Claude Code 提前到第一批。

\subsection{数据模型}

\noindent\begin{minipage}{\linewidth}
\centering
\captionof{table}{SQLite 核心表}
\small
\begin{tabularx}{\linewidth}{@{}p{3.2cm} X@{}}
\toprule
\textbf{表} & \textbf{关键字段} \\
\midrule
\code{projects} & id, name, path, default\_branch, created\_at \\
\code{assistants} & id, name, avatar, persona, proactivity, created\_at \\
\code{tasks} & id, project\_id, parent\_id, title, goal, status, assignee\_id, budget\_json, created\_at, updated\_at \\
\code{task\_edges} & from\_task\_id, to\_task\_id, type \\
\code{workers} & id, task\_id, role, runtime, model, worktree\_path, session\_ref, status, cost \\
\code{runs} & id, worker\_id, attempt, started\_at, ended\_at, result, error \\
\code{events} & id, task\_id, worker\_id, type, payload\_json, created\_at（append-only） \\
\code{artifacts} & id, task\_id, type, path, hash, created\_at \\
\code{memory\_entries} & id, scope, project\_id, type, content, source\_task\_id, confidence, expires\_at \\
\code{approvals} & id, task\_id, type, requested\_by, status, decided\_by, decided\_at \\
\code{budgets} & id, scope, limit\_json, used\_json, period \\
\code{model\_calls} & id, worker\_id, model, tokens\_in, tokens\_out, cost, created\_at \\
\bottomrule
\end{tabularx}
\end{minipage}

\subsection{仓库结构}

\begin{lstlisting}[caption={OpenCoxswain 仓库结构建议},label={lst:repo}]
opencoxswain/
├─ apps/
│  ├─ desktop/          # Tauri 2 + React
│  └─ daemon/           # Node/TS sidecar
├─ packages/
│  ├─ core/             # 任务契约、状态机、调度器
│  ├─ adapters/         # opencode、claude-code、codex
│  ├─ memory/           # Markdown + FTS5
│  └─ shared/           # 类型、事件、工具
├─ docs/
│  ├─ prd/              # 本文件
│  ├─ architecture/     # 架构与 ADR
│  └─ adr/              # 决策记录
├─ scripts/
├─ .github/workflows/
├─ LICENSE
├─ README.md
├─ CONTRIBUTING.md
└─ SECURITY.md
\end{lstlisting}

\subsection{关键流程}

\begin{enumerate}
  \item \textbf{Intake}：用户对总管说目标；总管读全局与项目记忆，最多问 1--3 个阻塞性问题。
  \item \textbf{Plan}：生成任务单；拆成 2--5 个子任务，形成 DAG；计算预算与风险。
  \item \textbf{Approve}：向用户展示执行清单；未确认不 spawn 员工。
  \item \textbf{Dispatch}：为每个任务创建 worktree、员工实例、会话；按依赖调度。
  \item \textbf{Execute}：员工在沙箱内执行；事件流回 Daemon；日志本地留存。
  \item \textbf{Review}：校验员独立复跑验收标准；输出通过/失败/未验证。
  \item \textbf{Summarize}：总管按三块格式汇总；用户审批后合并、推送或返工。
  \item \textbf{Remember}：验收通过后提取候选记忆；用户确认后写入。
\end{enumerate}

% ============================================================
%  11. 非功能需求
% ============================================================
\needspace{9cm}
\section{非功能需求}

\noindent\begin{minipage}{\linewidth}
\centering
\captionof{table}{非功能需求}
\small
\begin{tabularx}{\linewidth}{@{}p{2.6cm} X@{}}
\toprule
\textbf{维度} & \textbf{要求} \\
\midrule
响应 & UI 事件小于 200ms；任务状态更新小于 1s。 \\
可靠性 & 崩溃后从持久化状态恢复；未完成任务不丢。 \\
可回放 & 每个任务有完整事件时间线，可导出。 \\
安全 & 密钥不进 prompt；危险动作人审；审计日志不可篡改。 \\
隐私 & 默认本地；遥测 opt-in；记忆可一键导出/删除。 \\
跨平台 & v0 只做 Windows 10/11；架构上为 macOS/Linux 留接口。 \\
可扩展 & runtime adapter、角色模板、任务模板可插件化。 \\
成本 & 每任务/每日预算硬上限；模型分层可配置。 \\
\bottomrule
\end{tabularx}
\end{minipage}

% ============================================================
%  12. 版本规划与验收
% ============================================================
\section{版本规划与验收}
\label{sec:roadmap}

\subsection{v0：跑通闭环（不做记忆、不做 IM）}

\textbf{范围}：

\begin{itemize}
  \item 单项目、单总管、2 个员工（实施 + 校验）。
  \item 结构化任务单、执行清单确认、DAG 拆解（2--4 个）。
  \item git worktree 隔离、状态机、极简看板。
  \item 预算上限、危险动作人审、事件日志。
  \item 汇总三块。
\end{itemize}

\begin{keybox}[v0 验收剧本]
在一个真实仓库里说：“修复设置页夜间模式切换后不生效的问题。”

\begin{enumerate}
  \item 总管 1 分钟内产出任务单，字段完整。
  \item 用户确认后，spawn 2 个员工：实施员在独立 worktree 修改并跑测试；校验员独立复跑并出报告。
  \item 看板实时显示两人状态；用户可暂停其中一个。
  \item 总管按三块格式汇报；用户查看 diff 后点通过。
  \item 全程可回放；总花费小于 \$2；无人工接管代码。
\end{enumerate}

\textbf{通过标准}：上述 5 条全部满足，且失败时能正确进入 blocked/重试流程。
\end{keybox}

\subsection{v1：记忆 + 完整工作台}

\textbf{范围}：

\begin{itemize}
  \item 双层记忆（全局 + 项目）、写入门槛、混合检索、冲突处理。
  \item 记忆检查器、待决策收件箱、DAG 可视化、产物 diff 审阅。
  \item 任务模板、失败重试/升级、成本报表。
  \item 首次使用引导、项目记忆草稿。
\end{itemize}

\textbf{验收}：开一个新项目，不用重讲技术栈和偏好，总管能直接给出符合习惯的任务单；并能主动指出“这条旧约束不能硬搬”。

\subsection{v2：多项目 + 远程 + 自动化}

\textbf{范围}：

\begin{itemize}
  \item 全局路由、多项目记忆、IM 入口（飞书/Telegram）、定时/触发任务。
  \item 云端执行体、团队共享记忆（可选）、插件/技能市场（可选）。
\end{itemize}

\textbf{验收}：在手机上给总管发一句“把昨天那个 bug 的回归测试补了”，它能在云端或本地执行，并把待决策项推回 IM。

\subsection{Windows 专项验收}

\noindent\begin{minipage}{\linewidth}
\centering
\captionof{table}{Windows 专项验收清单}
\small
\begin{tabularx}{\linewidth}{@{}p{3.6cm} X@{}}
\toprule
\textbf{项目} & \textbf{验收标准} \\
\midrule
安装/卸载 & Windows 10 与 Windows 11 上均可安装、卸载；不残留后台进程。 \\
WebView2 & 检测运行时；缺失时引导安装或自带 bootstrapper。 \\
密钥存储 & API Key 存 Credential Manager/DPAPI；日志与崩溃报告中不出现明文。 \\
中文/空格路径 & 仓库路径包含中文、空格时，worktree、日志、产物路径均正常。 \\
超长路径 & 超过 260 字符的路径不崩溃；必要时启用长路径支持或给出明确错误。 \\
git worktree & 创建、切换、删除、异常清理均正常；文件锁有重试。 \\
杀软误报 & Defender 扫描通过；发布前在多引擎上抽查。 \\
自动更新 & 从 GitHub Releases 下载、校验、安装、重启；失败可回滚。 \\
进程清理 & 退出 App 时清理 sidecar、OpenCode 子进程、临时端口。 \\
崩溃恢复 & 强杀进程后重启，未完成任务从持久化状态恢复。 \\
\bottomrule
\end{tabularx}
\end{minipage}

\subsection{里程碑}

\noindent\begin{minipage}{\linewidth}
\centering
\captionof{table}{里程碑粗估（1 人全职 + AI 辅助）}
\small
\begin{tabularx}{\linewidth}{@{}p{1.6cm} X p{2.6cm}@{}}
\toprule
\textbf{阶段} & \textbf{内容} & \textbf{粗估} \\
\midrule
M0 & 方案确认、技术选型、OpenCode Windows spike、原型验证。 & 1 周 \\
M1 & v0 闭环。 & 4--6 周 \\
M2 & 自用内测、修体验。 & 2 周 \\
M3 & v1（记忆 + 完整工作台）。 & +4--6 周 \\
M4 & v2（多项目 + IM + 云端）。 & +6--8 周 \\
\bottomrule
\end{tabularx}
\end{minipage}

% ============================================================
%  13. 成功指标
% ============================================================
\section{成功指标}

\subsection*{北极星指标}

\begin{keybox}[北极星]
每周由总管端到端完成、并经用户验收通过的任务数。
\end{keybox}

\subsection*{辅助指标}

\noindent\begin{minipage}{\linewidth}
\centering
\captionof{table}{v0 目标指标}
\small
\begin{tabularx}{\linewidth}{@{}p{6.0cm} X@{}}
\toprule
\textbf{指标} & \textbf{v0 目标} \\
\midrule
任务一次验收通过率 & 大于等于 60\%。 \\
单任务人工干预次数 & 小于等于 2。 \\
从下达到首个产物时间 & 中位数小于 10 分钟。 \\
单任务平均成本 & 有上限、可解释。 \\
任务单被用户手改比例 & 小于 30\%（太高说明拆解差）。 \\
全链路可回放覆盖率 & 100\%。 \\
\bottomrule
\end{tabularx}
\end{minipage}

\subsection*{反指标}

\begin{itemize}
  \item 简单任务被错误路由到多员工的比例。
  \item token 浪费率（失败重试/无效探索占比）。
  \item 用户因不信任而频繁接管的比例。
\end{itemize}

% ============================================================
%  14. 风险与对策
% ============================================================
\section{风险与对策}

\noindent\begin{minipage}{\linewidth}
\centering
\captionof{table}{主要风险与对策}
\small
\begin{tabularx}{\linewidth}{@{}p{3.4cm} p{4.4cm} X@{}}
\toprule
\textbf{风险} & \textbf{影响} & \textbf{对策} \\
\midrule
拆解质量差 & 全队跑偏 & 强模型 + 模板 + 用户确认 + 小步验证。 \\
并行任务不独立 & 文件冲突、返工 & DAG 依赖 + 写域声明 + 冲突转人审。 \\
成本失控 & 用户流失 & 预算熔断 + 并发上限 + 模型分层 + 简单任务单员工。 \\
权限/安全 & 数据损坏、密钥泄露 & 白名单 + 沙箱 + 人审 + 审计 + 脱敏。 \\
记忆污染 & 越用越蠢 & 写入门槛 + 来源标注 + 用户可编辑 + 过期。 \\
平台依赖 & runtime 升级即坏 & adapter 抽象 + 版本兼容测试 + 降级方案。 \\
用户不信任 & 不敢放权 & 全链路可观测 + 可回放 + 随时接管。 \\
产品过重 & 简单任务也慢 & 轻路径：单员工直接干，不走完整编排。 \\
OpenCode Windows 不稳 & v0 卡住 & M0 spike；准备 CLI 包装与 Claude Code 备选。 \\
\bottomrule
\end{tabularx}
\end{minipage}

% ============================================================
%  15. 非目标
% ============================================================
\section{非目标（v0--v1 明确不做）}

\begin{itemize}
  \item 通用自主 Agent（什么都能干）。
  \item 非编码场景的一站式自动化。
  \item 多人实时协作、企业权限体系。
  \item 自研大模型。
  \item 移动端原生 App。
  \item 插件市场/技能商店。
  \item 云端多租户 SaaS。
\end{itemize}

% ============================================================
%  16. 开源策略
% ============================================================
\needspace{10cm}
\section{开源策略}
\label{sec:opensource}

\subsection{许可证}

\noindent\begin{minipage}{\linewidth}
\centering
\captionof{table}{Apache-2.0 与 AGPL-3.0 对比}
\small
\begin{tabularx}{\linewidth}{@{}p{2.6cm} X X@{}}
\toprule
\textbf{维度} & \textbf{Apache-2.0（建议）} & \textbf{AGPL-3.0（备选）} \\
\midrule
采用门槛 & 低，公司也愿意用。 & 高，部分公司直接禁用。 \\
专利授权 & 明确授予专利许可。 & 有专利条款，但更复杂。 \\
copyleft & 无。 & 强 copyleft；网络服务也需开源。 \\
适合目标 & 先做采用、做社区、做标准。 & 防止别人拿去做闭源云服务。 \\
\bottomrule
\end{tabularx}
\end{minipage}

\textbf{建议}：Apache-2.0。当前不做商业化，优先降低采用门槛；如果未来出现“别人直接托管收费”的风险，再考虑 AGPL 或双许可证。

\subsection{仓库与 CI}

单仓库 monorepo，结构见第 \ref{sec:architecture} 节。CI 跑 \code{windows-latest}：

\begin{itemize}
  \item \code{pnpm install} + \code{pnpm lint} + \code{pnpm test}。
  \item \code{cargo fmt --check} + \code{cargo clippy}（Tauri 壳）。
  \item \code{pnpm tauri build} 产物打包。
  \item 发布：推 \code{v*} 标签 $\rightarrow$ 构建 $\rightarrow$ 签名（可选）$\rightarrow$ GitHub Release $\rightarrow$ updater manifest。
\end{itemize}

\subsection{分发与代码签名}

\begin{itemize}
  \item \textbf{GitHub Releases}：主渠道；安装包 + 校验和 + 更新清单。
  \item \textbf{winget}：Windows 用户最顺的安装方式；稳定后提交 manifest。
  \item \textbf{Microsoft Store}：以后再说；需要额外的打包与合规工作。
  \item \textbf{代码签名}：不签名会触发 SmartScreen“未知发布者”。开源可申请 SignPath.io 免费签名，或使用 Azure Trusted Signing；自用阶段可以先不签。
\end{itemize}

\subsection{社区文档}

\begin{itemize}
  \item \code{README.md}：一句话、截图、安装、快速开始、架构图。
  \item \code{CONTRIBUTING.md}：开发环境、代码规范、提交规范、good first issue。
  \item \code{SECURITY.md}：漏洞报告渠道、支持版本、响应时间。
  \item \code{CODE\_OF\_CONDUCT.md}：行为准则。
  \item \code{ROADMAP.md}：v0/v1/v2 与当前进度。
  \item \code{docs/adr/}：关键决策记录（为什么 Windows 优先、为什么 OpenCode 首发、为什么不用向量库）。
\end{itemize}

\subsection{遥测与隐私}

默认零遥测；崩溃报告做成显式 opt-in；所有记忆、日志、产物默认存本地，用户可一键导出/删除。

\subsection{版本号}

SemVer；v1 之前保持 \code{0.x}；破坏性变更升 minor，功能升 minor，修复升 patch。具体规则在 \code{CONTRIBUTING.md} 里写清楚。

% ============================================================
%  17. 命名
% ============================================================
\section{命名：OpenCoxswain}
\label{sec:naming}

\subsection{最终命名}

\begin{keybox}[产品名]
\textbf{OpenCoxswain}（中文名：\textbf{舵手}）\\
The coxswain doesn't row. / 舵手不划桨。
\end{keybox}

\subsection{为什么不是 OpenSuperAgent}

在定名前，我们核查过 \code{OpenSuperAgent}，结论是\textbf{冲突严重，不建议使用}：

\noindent\begin{minipage}{\linewidth}
\centering
\captionof{table}{OpenSuperAgent 命名占用情况（2026-11-09 核查）}
\small
\begin{tabularx}{\linewidth}{@{}p{4.0cm} p{4.2cm} X@{}}
\toprule
\textbf{占用方} & \textbf{情况} & \textbf{影响} \\
\midrule
GitHub 组织 \code{OpenSuperAgent} & 已存在；2025-06 创建；AI 社区；1 个公开仓库。 & 拿不到同名组织，仓库只能挂个人账号或另起组织。 \\
\makecell[l]{\code{thefullgc/}\\\code{opensuperagent}} & 2026-11-02 创建；51 star；TypeScript；Apache-2.0；简介“DeepSeek Super-agent 开源落地实现”；npm \code{opensuperagent@0.1.1}。 & 同领域、同名字、已有社区和 npm 包，搜索流量会被它吃掉。 \\
\makecell[l]{\code{xinyuchen-pku/}\\\code{OpenSuperAgent}} & 2026-11-02 创建；Python；MIT；简介“Local-first repository agents with explicit permissions, isolated Git worktrees, and reviewable results.” & 概念几乎和 v0 一模一样：本地优先、权限、worktree 隔离、可审阅结果。 \\
\code{superagent} & 16.6k star 的知名 HTTP 客户端；npm \code{superagent@10.x}。 & \code{OpenSuperAgent} 会被读成“开源的 SuperAgent”，品牌混淆。 \\
npm \code{opensuperagent} & 已被占用。 & 发不了同名 CLI 包；只能用 scoped 包或改名。 \\
PyPI \code{opensuperagent} & 未占用。 & 但不足以抵消 GitHub/npm 的冲突。 \\
\bottomrule
\end{tabularx}
\end{minipage}

\textbf{结论}：名字太“品类化”（Open + Super + Agent），既难搜又难形成品牌；两个同名仓库还都在一周内出现，其中一个与 v0 概念高度重合。弃用。

\subsection{名字故事}

Coxswain 是赛艇上的舵手：他\textbf{不划桨}，负责掌舵、喊节奏、指挥桨手。桨手负责出力，舵手负责让整条船往同一个方向走。

这正是 OpenCoxswain 的产品哲学：

\begin{itemize}
  \item 用户是“定方向的人”：说目标、做决策、验收。
  \item 总管是“舵手”：不亲自划桨，负责拆解、派发、盯进度、汇总。
  \item 员工是“桨手”：各自独立上下文、独立 worktree，只对自己的任务负责。
\end{itemize}

“Open” 表示开源、开放、执行体可插拔；“Coxswain” 提供记忆点和故事。

\newpage
\subsection{命名体系}

\noindent\begin{minipage}{\linewidth}
\centering
\captionof{table}{OpenCoxswain 命名体系}
\small
\begin{tabularx}{\linewidth}{@{}p{3.2cm} X@{}}
\toprule
\textbf{项} & \textbf{名称} \\
\midrule
产品名 & OpenCoxswain \\
中文名 & 舵手 \\
桌面 App & OpenCoxswain \\
仓库 & \code{opencoxswain} \\
CLI & \code{opencoxswain} \\
npm scope & \code{@opencoxswain/core}、\code{@opencoxswain/adapters}、\code{@opencoxswain/memory}、\code{@opencoxswain/desktop} \\
英文 tagline & The coxswain doesn't row. \\
中文 tagline & 舵手不划桨。 \\
定位补充 & The open-source, local-first AI agent manager for Windows. \\
\bottomrule
\end{tabularx}
\end{minipage}

\subsection{README 首屏}

\begin{lstlisting}[caption={README 首屏文案建议},label={lst:readme},basicstyle=\ttfamily\scriptsize,aboveskip=2pt,belowskip=2pt]
# OpenCoxswain · 舵手

> The coxswain doesn't row.
> 舵手不划桨。

OpenCoxswain is a local-first, open-source AI agent manager for Windows.
You talk to one coxswain. It plans, dispatches, and reviews a crew of AI workers
while you make the decisions.

## Why

AI coding agents are getting better at doing the work.
They are still bad at managing the work.

OpenCoxswain gives you one persistent manager:
- it plans and splits tasks,
- it dispatches isolated workers,
- it reviews results,
- it remembers what worked,
- and it brings back only the decisions you need to make.

## Status

v0, early. Windows first. OpenCode first.
\end{lstlisting}

\subsection{定名前检查清单}

\begin{itemize}
  \item GitHub 组织名：\code{opencoxswain} 是否可注册（或使用个人账号 + 仓库名）。
  \item npm scope：\code{@opencoxswain} 是否可注册；备用 \code{opencoxswain-ai}。
  \item 域名：\code{opencoxswain.com} / \code{.dev} / \code{.app}。
  \item 商标：软件类（第 9 类、第 42 类）检索；至少确认没有同领域大厂。
  \item Windows Store 名称、GitHub Pages、社交媒体 handle。
  \item 搜索“OpenCoxswain”确认没有同领域项目。
\end{itemize}

\subsection{备选名字}

\noindent\begin{minipage}{\linewidth}
\centering
\captionof{table}{备选名字（如果 OpenCoxswain 遇到不可用问题）}
\small
\begin{tabularx}{\linewidth}{@{}p{2.8cm} p{1.8cm} X@{}}
\toprule
\textbf{英文} & \textbf{中文} & \textbf{一句话故事} \\
\midrule
Stationmaster & 总调度 & 站长不驾驶火车，只排班、控站台、保准点。 \\
Seneschal & 总管 & 中世纪庄园总管，替领主管理全部家仆。 \\
Impresario & 班主 & 班主不登台，组戏班、排演出、管艺人。 \\
Innkeeper & 掌柜 & 掌柜不跑堂，管伙计、管账、管货源。 \\
\bottomrule
\end{tabularx}
\end{minipage}

\subsection{已撞名、不建议使用}

\begin{itemize}
  \item \textbf{Showrunner}：已有一个几乎同类的“多 coding-agent 终端 UI”项目。
  \item \textbf{Maestro}：已有知名 UI 测试工具。
  \item \textbf{Conductor}：Netflix Conductor 已占。
  \item \textbf{Foreman}：Foreman 生命周期管理已占。
  \item \textbf{Bosun}：已有监控系统。
  \item \textbf{Gaffer}：已有 GafferHQ 视觉特效软件。
  \item \textbf{Steward / Majordomo / Butler}：通用词，且已有大量同名项目。
  \item \textbf{CrewAI}：直接撞 LangChain 生态的 CrewAI。
\end{itemize}

% ============================================================
%  18. 待拍板与下一步
% ============================================================
\section{待拍板与下一步}

\subsection{还需要拍板的问题}

\noindent\begin{minipage}{\linewidth}
\centering
\captionof{table}{待拍板问题}
\small
\begin{tabularx}{\linewidth}{@{}p{0.8cm} p{5.2cm} X@{}}
\toprule
\textbf{\#} & \textbf{问题} & \textbf{建议} \\
\midrule
1 & 许可证最终选 Apache-2.0 还是 AGPL-3.0？ & Apache-2.0；先降低采用门槛。 \\
2 & 是否随 App 捆绑 OpenCode，还是要求用户自行安装？ & v0 捆绑或自动下载；降低首次使用门槛。 \\
3 & Windows 最低版本？ & Windows 10 1809+ / Windows 11；覆盖大多数用户。 \\
4 & 是否购买代码签名证书？ & 公开分发前上 SignPath 或 Azure Trusted Signing。 \\
5 & 遥测默认关闭，崩溃报告是否 opt-in？ & 默认关闭；崩溃报告显式 opt-in。 \\
6 & 是否注册 \code{opencoxswain} 组织与域名？ & 定名后立即注册，避免被抢。 \\
\bottomrule
\end{tabularx}
\end{minipage}

\subsection{下一步}

\begin{enumerate}
  \item 冻结 v0 范围与验收剧本。
  \item 做 OpenCode Windows spike：验证 serve/SDK、worktree、权限、事件流。
  \item 写任务单 schema、状态机、RuntimeAdapter 接口的第一版代码。
  \item 搭一个薄薄的可运行纵切：一句需求 $\rightarrow$ 1 个总管 $\rightarrow$ 1 个员工 $\rightarrow$ 1 个 diff。
  \item 自用 2 周，记录每次“总管拆错/员工跑偏/验收失败”的案例，再决定 v1 记忆的优先级。
\end{enumerate}

% ============================================================
%  Appendix
% ============================================================
\appendix
\section{任务单 JSON Schema（简化版）}

\begin{lstlisting}[language=Java,caption={TaskContract JSON Schema（节选）},label={lst:schema}]
{
  "$schema": "https://json-schema.org/draft/2020-12/schema",
  "title": "TaskContract",
  "type": "object",
  "required": [
    "task_id", "project_id", "title", "goal",
    "deliverables", "acceptance_criteria", "budget"
  ],
  "properties": {
    "task_id": { "type": "string" },
    "parent_id": { "type": ["string", "null"] },
    "project_id": { "type": "string" },
    "title": { "type": "string" },
    "goal": { "type": "string" },
    "context_refs": { "type": "array", "items": { "type": "string" } },
    "deliverables": {
      "type": "array",
      "items": {
        "type": "object",
        "required": ["type"],
        "properties": {
          "type": { "enum": ["diff", "file", "report", "link"] },
          "scope": { "type": "string" }
        }
      }
    },
    "acceptance_criteria": {
      "type": "array",
      "items": {
        "type": "object",
        "required": ["id", "type", "desc"],
        "properties": {
          "id": { "type": "string" },
          "type": { "enum": ["command", "diff", "rubric", "human"] },
          "desc": { "type": "string" },
          "check": { "type": "string" }
        }
      }
    },
    "dependencies": { "type": "array", "items": { "type": "string" } },
    "budget": {
      "type": "object",
      "properties": {
        "tokens": { "type": "integer" },
        "usd": { "type": "number" },
        "minutes": { "type": "integer" }
      }
    }
  }
}
\end{lstlisting}

\section{RuntimeAdapter 接口（完整版）}

\begin{lstlisting}[language=Java,caption={RuntimeAdapter 完整接口},label={lst:adapter-full}]
export interface RuntimeAdapter {
  readonly name: string;
  readonly version: string;

  createSession(opts: SessionOptions): Promise<RuntimeSession>;
  send(session: RuntimeSession, task: TaskContract): AsyncIterable<RuntimeEvent>;
  cancel(session: RuntimeSession): Promise<void>;
  collect(session: RuntimeSession): Promise<WorkerResult>;
  dispose(session: RuntimeSession): Promise<void>;
  healthCheck(): Promise<HealthStatus>;
}

export type RuntimeEvent =
  | { type: "started"; at: number }
  | { type: "message"; at: number; text: string }
  | { type: "tool_call"; at: number; tool: string; input: unknown }
  | { type: "tool_result"; at: number; tool: string; output: unknown }
  | { type: "artifact"; at: number; artifact: ArtifactRef }
  | { type: "progress"; at: number; percent: number; note?: string }
  | { type: "error"; at: number; message: string; retryable: boolean }
  | { type: "finished"; at: number; result: WorkerResult };

export interface SessionOptions {
  taskId: string;
  worktree: string;
  role: RoleTemplate;
  model: ModelConfig;
  approvalPolicy: ApprovalPolicy;
  env: Record<string, string>;
}

export interface RoleTemplate {
  name: "researcher" | "implementer" | "verifier" | "writer";
  allowedTools: string[];
  readOnly: boolean;
  outputSchema: string;
}

export interface ModelConfig {
  provider: string;
  model: string;
  apiKeyRef: string;       // 只存引用，不存明文
  maxTokens: number;
  temperature?: number;
}

export interface ApprovalPolicy {
  requireHumanFor: Array<
    "delete" | "push" | "publish" | "pay" | "external_send"
  >;
  autoApproveUnderUsd?: number;
}
\end{lstlisting}

\section{SQLite 表结构（节选）}

\begin{lstlisting}[language=SQL,caption={SQLite 表结构（节选）},label={lst:sql}]
CREATE TABLE projects (
  id             TEXT PRIMARY KEY,
  name           TEXT NOT NULL,
  path           TEXT NOT NULL,
  default_branch TEXT NOT NULL DEFAULT 'main',
  created_at     INTEGER NOT NULL
);

CREATE TABLE tasks (
  id           TEXT PRIMARY KEY,
  project_id   TEXT NOT NULL REFERENCES projects(id),
  parent_id    TEXT REFERENCES tasks(id),
  title        TEXT NOT NULL,
  goal         TEXT NOT NULL,
  status       TEXT NOT NULL,
  assignee_id  TEXT,
  budget_json  TEXT NOT NULL,
  created_at   INTEGER NOT NULL,
  updated_at   INTEGER NOT NULL
);

CREATE TABLE task_edges (
  from_task_id TEXT NOT NULL REFERENCES tasks(id),
  to_task_id   TEXT NOT NULL REFERENCES tasks(id),
  type         TEXT NOT NULL DEFAULT 'depends_on',
  PRIMARY KEY (from_task_id, to_task_id)
);

CREATE TABLE events (
  id         INTEGER PRIMARY KEY AUTOINCREMENT,
  task_id    TEXT NOT NULL,
  worker_id  TEXT,
  type       TEXT NOT NULL,
  payload    TEXT NOT NULL,
  created_at INTEGER NOT NULL
);

CREATE TABLE memory_entries (
  id            TEXT PRIMARY KEY,
  scope         TEXT NOT NULL CHECK (scope IN ('global', 'project')),
  project_id    TEXT,
  type          TEXT NOT NULL,
  content       TEXT NOT NULL,
  source_task_id TEXT,
  confidence    REAL NOT NULL DEFAULT 0.5,
  expires_at    INTEGER,
  created_at    INTEGER NOT NULL
);

CREATE VIRTUAL TABLE memory_fts USING fts5(
  content, tags, content='memory_entries', content_rowid='rowid'
);
\end{lstlisting}

\needspace{8cm}
\section{命名占用核查记录}

\noindent\begin{minipage}{\linewidth}
\centering
\captionof{table}{命名核查记录（2026-11-09）}
\small
\begin{tabularx}{\linewidth}{@{}p{3.6cm} p{2.4cm} X@{}}
\toprule
\textbf{名称} & \textbf{状态} & \textbf{说明} \\
\midrule
\code{OpenSuperAgent} & \bad & GitHub 组织已存在；两个同名仓库；npm 包已占用；与 v0 概念高度重合。 \\
\code{OpenCoxswain} & \good & GitHub 搜索 0 结果；npm 未占用；PyPI 未占用。 \\
\code{opencoxswain} & \good & GitHub 搜索 0 结果；npm 未占用。 \\
\code{open-coxswain} & \good & GitHub 搜索 0 结果；npm 未占用。 \\
\code{coxswain-ai} & \warn & GitHub 仅 1 个 1 star 仓库；npm 未占用。 \\
\code{Coxswain} & \warn & 已有划船/编译器项目同名；作为产品名仍可用，但 npm 包名需加 scope。 \\
\bottomrule
\end{tabularx}
\end{minipage}

\begin{notebox}[核查边界]
以上为公开平台检索结果，不构成商标法律意见。正式发布前，建议在目标市场做一次软件类商标检索，并确认域名、组织名、npm scope 可用。
\end{notebox}

\end{document}
